problem:  
Let \(G\) be a connected Lie group with compact subgroups \(K \subset H \subset G\), and let \(M = G/K\) be a homogeneous Riemannian manifold equipped with a \(G\)-invariant metric \(g\).  
Assume that \((M,g)\) is *weakly symmetric* in the sense of Selberg–Vinberg.  
Denote by \(\pi : G/K \to G/H\) the canonical projection.  

Suppose there exists  
- a nontrivial smooth map \(f: G/K \to G/K\),  
- a Lie group automorphism \(\varphi: G \to G\) preserving both \(H\) and \(K\),  
and an induced algebraic automorphism \(\psi : G/H \to G/H\) satisfying  
\[
\psi(gH) = \varphi(g)H,
\]
such that the following *twisted \(G\)-equivariance* and diagrammatic relations hold:
\[
f(g \cdot x) = \varphi(g) \cdot f(x), \qquad
\begin{CD}
G/K @>f>> G/K \\
@V{\pi}VV  @VV{\pi}V \\
G/H @>>\psi> G/H
\end{CD}
\]
and \(f\) is a *projective transformation* of \((M,g)\), i.e., \(f\) sends unparameterized geodesics to unparameterized geodesics.  

**Question:**  
If \((M,g)\) is weakly symmetric and such a nontrivial \(f\) (not induced by an element of \(G\)) exists, must \((M,g)\) be locally isometric to a globally symmetric space?  
Equivalently, can a *non-symmetric weakly symmetric* space admit a twisted \(G\)-equivariant projective self-map that fits a commutative diagram over \(G/H\)?  

Why is it a "good" problem:  
This formulation precisely captures an open geometric rigidity phenomenon at the intersection of homogeneous differential geometry, Lie group theory, and projective differential geometry. It asks whether an extra layer of symmetry—expressed by a *twisted equivariant diagrammatic projective self-map*—forces a weakly symmetric manifold to collapse to a fully symmetric one.  

The problem probes how categorical self-consistency (the commutative diagram condition) and dynamical symmetry (geodesic preservation) interact in homogeneous spaces. It generalizes classical rigidity questions (like the Lichnerowicz–Obata theorem for projective transformations) into the weakly symmetric regime, where the standard tools of symmetric-space theory no longer directly apply.  

The condition \(f(gx) = \varphi(g)f(x)\) introduces a genuinely new feature absent from prior rigidity literature: a coupling between group automorphisms and geometric projective transformations. Whether such an intertwining can exist without full symmetry is completely unknown and could reveal intrinsic constraints on the relationship between weak symmetry, automorphism structure, and curvature behavior.  

Thus, the problem offers a conceptually simple geometric setup with potentially far-reaching consequences and cross-field depth—an archetypal “good” modern research problem.